\chapter{Conversion between Sanskrit and Pāli forms}\label{chap:convsktpali}

This part is taken from Wikipedia\footnote{\url{en.wikipedia.org/wiki/Pāli}} to help Pāli learners be familiar with the way Sanskrit words transform into Pāli. Knowing this can enhance ability to find meaning of words not described in existing Pāli dictionaries.\footnote{If you take the correspondences between Pāli and Sanskrit seriously, a recommended reading is Thomas Oberlies's \emph{Pāli: A Grammar of the language of the Theravāda Tipiṭaka} (2001, Walter de Gruyter).}

\section{Vowels and diphthongs}

\begin{itemize}
	\item Sanskrit \pali{ai} and \pali{au} always monophthongize to Pāli \pali{e} and \pali{o}, respectively
		\begin{itemize}
			\item \pali{maitrī} (friendliness, benevolence) → \pali{mettā}
			\item \pali{auṣadha} (medical herb) → \pali{osadha}
		\end{itemize}
	\item Sanskrit \pali{āya}, \pali{ayā} and \pali{avā} reduce to Pāli \pali{ā}
		\begin{itemize}
			\item \pali{katipayāha} (someone) → \pali{katipāha}
			\item \pali{vaihāyasa} (sky-dwelling) → \pali{vehāsa}
			\item \pali{yāvagū} (barley) → \pali{yāgu}
		\end{itemize}
	\item Sanskrit \pali{aya} and \pali{ava} likewise often reduce to Pāli \pali{e} and \pali{o}
		\begin{itemize}
			\item \pali{dhārayati} (one maintains, one holds) → \pali{dhāreti}
			\item \pali{avatāra} (descent) → \pali{otāra}
			\item \pali{bhavati} (one becomes) → \pali{hoti}
		\end{itemize}
	\item Sanskrit \pali{avi} and \pali{ayū} becomes Pāli \pali{e} (i.e. \pali{avi → ai → e}) and \pali{o}
		\begin{itemize}
			\item \pali{sthavira} (broad, thick, compact) → \pali{thera}
			\item \pali{mayūra} (peacock) → \pali{mora}
		\end{itemize}
	\item Sanskrit \pali{ṛ} appears in Pāli as \pali{a}, \pali{i} or \pali{u}, often agreeing with the vowel in the following syllable. \pali{ṛ} also sometimes becomes \pali{u} after labial consonants.
		\begin{itemize}
			\item \pali{kṛta} (done) → \pali{kata}
			\item \pali{tṛṣṇa} (thirst) → \pali{taṇha}
			\item \pali{smṛti} (remembrance, reminiscence) → \pali{sati}
			\item \pali{ṛṣi} (cleric) → \pali{isi}
			\item \pali{dṛṣṭi} (vision, sight) → \pali{diṭṭhi}
			\item \pali{ṛddhi} (growth, increase) → \pali{iddhi}
			\item \pali{ṛju} (straight) → \pali{uju}
			\item \pali{spṛṣṭa} (touched) → \pali{phuṭṭha}
			\item \pali{vṛddha} (old) → \pali{vuddha}
		\end{itemize}
	\item Sanskrit long vowels are shortened before a sequence of two following consonants.
		\begin{itemize}
			\item \pali{kṣānti} (patience, forbearance, endurance, indulgence) → \pali{khanti}
			\item \pali{rājya} (kingdom) → \pali{rajja}
			\item \pali{īśvara} (lord) → \pali{issara}
			\item \pali{tīrṇa} (crossed, surpassed) → \pali{tiṇṇa}
			\item \pali{pūrva} (east) → \pali{pubba}
		\end{itemize}
\end{itemize}

\section{Consonants}

\subsection{Sound changes}

\begin{itemize}
	\item The Sanskrit sibilants \pali{ś, ṣ,} and \pali{s} merge as Pāli \pali{s}
		\begin{itemize}
			\item \pali{śaraṇa} (protector, defender) → \pali{saraṇa}
			\item \pali{doṣa} (night, darkness) → \pali{dosa}
		\end{itemize}
	\item The Sanskrit stops \pali{ḍ} and \pali{ḍh} become \pali{ḷ} and \pali{ḷh} between vowels (as in Vedic)
		\begin{itemize}
			\item \pali{cakravāḍa} (cyclic) → \pali{cakkavāḷa}
			\item \pali{virūḍha} (mounted, sprouted) → \pali{virūḷha}
		\end{itemize}
\end{itemize}

\subsection{Assimilations}

\subsection*{General rules}

\begin{itemize}
	\item Many assimilations of one consonant to a neighboring consonant occurred in the development of Pāli, producing a large number of geminate (double) consonants. Since aspiration of a geminate consonant is only phonetically detectable on the last consonant of a cluster, geminate \pali{kh, gh, ch, jh, ṭh, ḍh, th, dh, ph} and \pali{bh} appear as \pali{kkh, ggh, cch, jjh, ṭṭh, ḍḍh, tth, ddh, pph} and \pali{bbh}, not as \pali{khkh, ghgh} etc.
	\item Initial consonant clusters are simplified to a single consonant.
		\begin{itemize}
			\item \pali{prāṇa} (respiration) → \pali{pāṇa} (not \pali{ppāṇa})
			\item \pali{sthavira} (compact, dense) → \pali{thera} (not \pali{tthera})
			\item \pali{dhyāna} (meditation) → \pali{jhāna} (not \pali{jjhāna})
			\item \pali{jñāti} (intelligence) → \pali{ñāti} (not \pali{ññāti})
		\end{itemize}
	\item When assimilation would produce a sequence of three consonants in the middle of a word, geminates are simplified until there are only two consonants in sequence.
		\begin{itemize}
			\item \pali{uttrāsa} (fear, terror) → \pali{uttāsa} (not \pali{utttāsa})
			\item \pali{mantra} (instrument of thought, speech) → \pali{manta} (not \pali{mantta})
			\item \pali{indra} (conqueror) → \pali{inda} (not \pali{indda})
			\item \pali{vandhya} (barren, fruitless, deprived) → \pali{vañjha} (not \pali{vañjjha})
		\end{itemize}
	\item The sequence \pali{vv} resulting from assimilation changes to \pali{bb}.
		\begin{itemize}
			\item \pali{sarva} (all, every, various) → \pali{savva} → \pali{sabba}
			\item \pali{pravrajati} (one moves forth) → \pali{pavvajati} → \pali{pabbajati}
			\item \pali{divya} (supernatural, wonderful, magical) → \pali{divva} → \pali{dibba}
			\item \pali{nirvāṇa} (deceased, extinguished; extinction, cessation, vanishing, disappearance) → \pali{nivvāṇa} → \pali{nibbāna}
		\end{itemize}
\end{itemize}

\subsection*{Total assimilation}

Total assimilation, where one sound becomes identical to a neighboring sound, is of two types: progressive, where the assimilated sound becomes identical to the following sound; and regressive, where it becomes identical to the preceding sound.

\subsection*{Regressive assimilations}

\begin{itemize}
	\item Internal visarga assimilates to a following voiceless stop or sibilant
		\begin{itemize}
			\item \pali{duḥkṛta} (=\pali{duṣkṛta}, wrong-done) → \pali{dukkata}
			\item \pali{duḥkha} (difficult, unagreeable) → \pali{dukkha}
			\item \pali{duḥprajña} (misknowledge) → \pali{duppañña}
			\item \pali{niḥkrodha} (=\pali{niṣkrodha}, wrath) → \pali{nikkodha}
			\item \pali{niḥpakva} (=\pali{niṣpakva}, well-cooked, decocted, infused) → \pali{nippakka}
			\item \pali{niḥśoka} (ugly, unhappy, inglorious)→ \pali{nissoka}
			\item \pali{niḥsattva} → \pali{nissatta}
		\end{itemize}
	\item In a sequence of two dissimilar Sanskrit stops, the first stop assimilates to the second stop
		\begin{itemize}
			\item \pali{vimukti} → \pali{vimutti}
			\item \pali{dugdha} → \pali{duddha}
			\item \pali{utpāda} → \pali{uppāda}
			\item \pali{pudgala} → \pali{puggala}
			\item \pali{udghoṣa} → \pali{ugghosa}
			\item \pali{adbhuta} → \pali{abbhuta}
			\item \pali{śabda} → \pali{sadda}
		\end{itemize}
	\item In a sequence of two dissimilar nasals, the first nasal assimilates to the second nasal
		\begin{itemize}
			\item \pali{unmatta} → \pali{ummatta}
			\item \pali{pradyumna} → \pali{pajjunna}
		\end{itemize}
	\item \pali{j} assimilates to a following \pali{ñ} (i.e., \pali{jñ} becomes \pali{ññ})
		\begin{itemize}
			\item \pali{prajñā} → \pali{paññā}
			\item \pali{jñāti} → \pali{ñāti}
		\end{itemize}
	\item The Sanskrit liquid consonants \pali{r} and \pali{l} assimilate to a following stop, nasal, sibilant, or \pali{v}
		\begin{itemize}
			\item \pali{mārga} → \pali{magga}
			\item \pali{karma} → \pali{kamma}
			\item \pali{varṣa} → \pali{vassa}
			\item \pali{kalpa} → \pali{kappa}
			\item \pali{sarva} → \pali{savva} → \pali{sabba}
		\end{itemize}
	\item \pali{r} assimilates to a following \pali{l}
		\begin{itemize}
			\item \pali{durlabha} → \pali{dullabha}
			\item \pali{nirlopa} → \pali{nillopa}
		\end{itemize}
	\item \pali{d} sometimes assimilates to a following \pali{v}, producing \pali{vv} → \pali{bb}
		\begin{itemize}
			\item \pali{udvigna} → \pali{uvvigga} → \pali{ubbigga}
			\item \pali{dvādaśa} → \pali{bārasa} (beside \pali{dvādasa})
		\end{itemize}
	\item \pali{t} and \pali{d} may assimilate to a following \pali{s} or \pali{y} when a morpheme boundary intervenes
		\begin{itemize}
			\item \pali{ut+sava} → \pali{ussava}
			\item \pali{ud+yāna} → \pali{uyyāna}
		\end{itemize}
\end{itemize}

\subsection*{Progressive assimilations}

\begin{itemize}
	\item Nasals sometimes assimilate to a preceding stop (in other cases epenthesis occurs)
		\begin{itemize}
			\item \pali{agni} (fire) → \pali{aggi}
			\item \pali{ātman} (self) → \pali{atta}
			\item \pali{prāpnoti} → \pali{pappoti}
			\item \pali{śaknoti} → \pali{sakkoti}
		\end{itemize}
	\item \pali{m} assimilates to an initial sibilant
		\begin{itemize}
			\item \pali{smarati} → \pali{sarati}
			\item \pali{smṛti} → \pali{sati}
		\end{itemize}
	\item Nasals assimilate to a preceding stop+sibilant cluster, which then develops in the same way as such clusters without following nasals
		\begin{itemize}
			\item \pali{tīkṣṇa} → \pali{tikṣa} → \pali{tikkha}
			\item \pali{lakṣmī} → \pali{lakṣī} → \pali{lakkhī}
		\end{itemize}
	\item The Sanskrit liquid consonants \pali{r} and \pali{l} assimilate to a preceding stop, nasal, sibilant, or \pali{v}
		\begin{itemize}
			\item \pali{prāṇa} → \pali{pāṇa}
			\item \pali{grāma} → \pali{gāma}
			\item \pali{śrāvaka} → \pali{sāvaka}
			\item \pali{agra} → \pali{agga}
			\item \pali{indra} → \pali{inda}
			\item \pali{pravrajati} → \pali{pavvajati} → \pali{pabbajati}
			\item \pali{aśru} → \pali{assu}
		\end{itemize}
	\item \pali{y} assimilates to preceding non-dental/retroflex stops or nasals
		\begin{itemize}
			\item \pali{cyavati} → \pali{cavati}
			\item \pali{jyotiṣ} → \pali{joti}
			\item \pali{rājya} → \pali{rajja}
			\item \pali{matsya} → \pali{macchya} → \pali{maccha}
			\item \pali{lapsyate} → \pali{lacchyate} → \pali{lacchati}
			\item \pali{abhyāgata} → \pali{abbhāgata}
			\item \pali{ākhyāti} → \pali{akkhāti}
			\item \pali{saṁkhyā} → \pali{saṅkhā} (but also \pali{saṅkhyā})
			\item \pali{ramya} → \pali{ramma}
		\end{itemize}
	\item \pali{y} assimilates to preceding non-initial \pali{v}, producing \pali{vv} → \pali{bb}
		\begin{itemize}
			\item \pali{divya} → \pali{divva} → \pali{dibba}
			\item \pali{veditavya} → \pali{veditavva} → \pali{veditabba}
			\item \pali{bhāvya} → \pali{bhavva} → \pali{bhabba}
		\end{itemize}
	\item \pali{y} and \pali{v} assimilate to any preceding sibilant, producing \pali{ss}
		\begin{itemize}
			\item \pali{paśyati} → \pali{passati}
			\item \pali{śyena} → \pali{sena}
			\item \pali{aśva} → \pali{assa}
			\item \pali{īśvara} → \pali{issara}
			\item \pali{kariṣyati} → \pali{karissati}
			\item \pali{tasya} → \pali{tassa}
			\item \pali{svāmin} → \pali{sāmī}
		\end{itemize}
	\item \pali{v} sometimes assimilates to a preceding stop
		\begin{itemize}
			\item \pali{pakva} → \pali{pakka}
			\item \pali{catvāri} → \pali{cattāri}
			\item \pali{sattva} → \pali{satta}
			\item \pali{dhvaja} → \pali{dhaja}
		\end{itemize}
\end{itemize}

\subsection*{Partial and mutual assimilation}

\begin{itemize}
	\item Sanskrit sibilants before a stop assimilate to that stop, and if that stop is not already aspirated, it becomes aspirated; e.g. \pali{śc, st, ṣṭ} and \pali{sp} become \pali{cch, tth, ṭṭh} and \pali{pph}
		\begin{itemize}
			\item \pali{paścāt} → \pali{pacchā}
			\item \pali{asti} → \pali{atthi}
			\item \pali{stava} → \pali{thava}
			\item \pali{śreṣṭha} → \pali{seṭṭha}
			\item \pali{aṣṭa} → \pali{aṭṭha}
			\item \pali{sparśa} → \pali{phassa}
		\end{itemize}
	\item In sibilant-stop-liquid sequences, the liquid is assimilated to the preceding consonant, and the cluster behaves like sibilant-stop sequences; e.g. \pali{str} and \pali{ṣṭr} become \pali{tth} and \pali{ṭṭh}
		\begin{itemize}
			\item \pali{śāstra} → \pali{śasta} → \pali{sattha}
			\item \pali{rāṣṭra} → \pali{raṣṭa} → \pali{raṭṭha}
		\end{itemize}
	\item \pali{t} and \pali{p} become \pali{c} before \pali{s}, and the sibilant assimilates to the preceding sound as an aspirate (i.e., the sequences \pali{ts} and \pali{ps} become \pali{cch})
		\begin{itemize}
			\item \pali{vatsa} → \pali{vaccha}
			\item \pali{apsaras} → \pali{accharā}
		\end{itemize}
	\item A sibilant assimilates to a preceding \pali{k} as an aspirate (i.e., the sequence \pali{kṣ} becomes \pali{kkh})
		\begin{itemize}
			\item \pali{bhikṣu} → \pali{bhikkhu}
			\item \pali{kṣānti} → \pali{khanti}
		\end{itemize}
	\item Any dental or retroflex stop or nasal followed by \pali{y} converts to the corresponding palatal sound, and the \pali{y} assimilates to this new consonant, i.e. \pali{ty, thy, dy, dhy, ny} become \pali{cc, cch, jj, jjh, ññ}; likewise \pali{ṇy} becomes \pali{ññ}. Nasals preceding a stop that becomes palatal share this change.
		\begin{itemize}
			\item \pali{tyajati} → \pali{cyajati} → \pali{cajati}
			\item \pali{satya} → \pali{sacya} → \pali{sacca}
			\item \pali{mithyā} → \pali{michyā} → \pali{micchā}
			\item \pali{vidyā} → \pali{vijyā} → \pali{vijjā}
			\item \pali{madhya} → \pali{majhya} → \pali{majjha}
			\item \pali{anya} → \pali{añya} → \pali{añña}
			\item \pali{puṇya} → \pali{puñya} → \pali{puñña}
			\item \pali{vandhya} → \pali{vañjhya} → \pali{vañjjha} → \pali{vañjha}
		\end{itemize}
	\item The sequence \pali{mr} becomes \pali{mb}, via the epenthesis of a stop between the nasal and liquid, followed by assimilation of the liquid to the stop and subsequent simplification of the resulting geminate.
		\begin{itemize}
			\item \pali{āmra} → \pali{ambra} → \pali{amba}
			\item \pali{tāmra} → \pali{tamba}
		\end{itemize}
\end{itemize}

\subsection{Epenthesis}

An epenthetic vowel is sometimes inserted between certain consonant-sequences. As with \pali{ṛ}, the vowel may be \pali{a, i}, or \pali{u}, depending on the influence of a neighboring consonant or of the vowel in the following syllable. \pali{i} is often found near \pali{i, y}, or palatal consonants; \pali{u} is found near \pali{u, v}, or labial consonants.

\begin{itemize}
	\item Sequences of stop + nasal are sometimes separated by \pali{a} or \pali{u}
		\begin{itemize}
			\item \pali{ratna} → \pali{ratana}
			\item \pali{padma} → \pali{paduma} (\pali{u} influenced by labial \pali{m})
		\end{itemize}
	\item The sequence \pali{sn} may become \pali{sin} initially
		\begin{itemize}
			\item \pali{snāna} → \pali{sināna}
			\item \pali{sneha} → \pali{sineha}
		\end{itemize}
	\item \pali{i} may be inserted between a consonant and \pali{l}
		\begin{itemize}
			\item \pali{kleśa} → \pali{kilesa}
			\item \pali{glāna} → \pali{gilāna}
			\item \pali{mlāyati} → \pali{milāyati}
			\item \pali{ślāghati} → \pali{silāghati}
		\end{itemize}
	\item An epenthetic vowel may be inserted between an initial sibilant and \pali{r}
		\begin{itemize}
			\item \pali{śrī} → \pali{sirī}
		\end{itemize}
	\item The sequence \pali{ry} generally becomes \pali{riy} (\pali{i} influenced by following \pali{y}), but is still treated as a two-consonant sequence for the purposes of vowel-shortening
		\begin{itemize}
			\item \pali{ārya} → \pali{arya} → \pali{ariya}
			\item \pali{sūrya} → \pali{surya} → \pali{suriya}
			\item \pali{vīrya} → \pali{virya} → \pali{viriya}
		\end{itemize}
	\item \pali{a} or \pali{i} is inserted between \pali{r} and \pali{h}
		\begin{itemize}
			\item \pali{arhati} → \pali{arahati}
			\item \pali{garhā} → \pali{garahā}
			\item \pali{barhiṣ} → \pali{barihisa}
		\end{itemize}
	\item There is sporadic epenthesis between other consonant sequences
		\begin{itemize}
			\item \pali{caitya} → \pali{cetiya} (not \pali{cecca})
			\item \pali{vajra} → \pali{vajira} (not \pali{vajja})
		\end{itemize}
\end{itemize}

\subsection{Other changes}

\begin{itemize}
	\item Any Sanskrit sibilant before a nasal becomes a sequence of nasal followed by \pali{h}, i.e. \pali{ṣṇ, sn}, and \pali{sm} become \pali{ṇh, nh}, and \pali{mh}
		\begin{itemize}
			\item \pali{tṛṣṇa} → \pali{taṇha}
			\item \pali{uṣṇīṣa} → \pali{uṇhīsa}
			\item \pali{asmi} → \pali{amhi}
	\end{itemize}
\item The sequence \pali{śn} becomes \pali{ñh}, due to assimilation of the \pali{n} to the preceding palatal sibilant
		\begin{itemize}
			\item \pali{praśna} → \pali{praśña} → \pali{pañha}
		\end{itemize}
	\item The sequences \pali{hy} and \pali{hv} undergo metathesis
		\begin{itemize}
			\item \pali{jihvā} → \pali{jivhā}
			\item \pali{gṛhya} → \pali{gayha}
			\item \pali{guhya} → \pali{guyha}
		\end{itemize}
	\item \pali{h} undergoes metathesis with a following nasal
		\begin{itemize}
			\item \pali{gṛhṇāti} → \pali{gaṇhāti}
		\end{itemize}
	\item \pali{y} is geminated between \pali{e} and a vowel
		\begin{itemize}
			\item \pali{śreyas} → \pali{seyya}
			\item \pali{Maitreya} → \pali{Metteyya}
		\end{itemize}
	\item Voiced aspirates such as \pali{bh} and \pali{gh} on rare occasions become \pali{h}
		\begin{itemize}
			\item \pali{bhavati} → \pali{hoti}
			\item \pali{-ebhiṣ} → \pali{-ehi}
			\item \pali{laghu} → \pali{lahu}
		\end{itemize}
	\item Dental and retroflex sounds sporadically change into one another
		\begin{itemize}
			\item \pali{jñāna} → \pali{ñāṇa} (not \pali{ñāna})
			\item \pali{dahati} → \pali{ḍahati} (beside Pāli \pali{dahati})
			\item \pali{nīḍa} → \pali{nīla} (not \pali{nīḷa})
			\item \pali{sthāna} → \pali{ṭhāna} (not \pali{thāna})
			\item \pali{duḥkṛta} → \pali{dukkaṭa} (beside Pāli \pali{dukkata})
			\item \pali{granthi}→ \pali{gaṇṭhi}
			\item \pali{pṛthivī} → \pali{paṭhavī/puṭhuvī} \\(beside Pāli \pali{pathavī/puthuvī/puthavī})
		\end{itemize}
\end{itemize}

\section{Exceptions}

There are several notable exceptions to the rules above; many of them are common Prakrit words rather than borrowings from Sanskrit.

\begin{itemize}
	\item \pali{ārya} (noble, pure) → \pali{ayya} (beside \pali{ariya})
	\item \pali{guru} (master) → \pali{garu} (adj.) (beside \pali{guru} (n.))
	\item \pali{puruṣa} (man) → \pali{purisa} (not \pali{purusa})
	\item \pali{vṛkṣa} (tree) → \pali{rukṣa} → \pali{rukkha} (not \pali{vukkha})
\end{itemize}

